\documentclass[11pt,a4paper]{article}

% ---------- Kodierung & Sprache ----------
\usepackage[T1]{fontenc}
\usepackage[utf8]{inputenc}
\usepackage[german]{babel}
\usepackage{lmodern}
\usepackage{microtype}
\usepackage{csquotes}

% ---------- Layout ----------
\usepackage[a4paper,margin=2.5cm]{geometry}
\usepackage{setspace}
\onehalfspacing
\usepackage{fancyhdr}
\pagestyle{fancy}
\fancyhf{}
\fancyhead[L]{\small Analyseprotokoll -- Q-Learning Lauf 1}
\fancyhead[R]{\small\thepage}
\renewcommand{\headrulewidth}{0.4pt}

% ---------- Tabellen ----------
\usepackage{booktabs}
\usepackage{tabularx}
\usepackage{multirow}
\usepackage{array}
\usepackage{caption}
\captionsetup{font=small,labelfont=bf}
\newcolumntype{L}[1]{>{\raggedright\arraybackslash}p{#1}}
\newcolumntype{C}[1]{>{\centering\arraybackslash}p{#1}}
\newcolumntype{R}[1]{>{\raggedleft\arraybackslash}p{#1}}

% ---------- Mathematik & Sonstiges ----------
\usepackage{amsmath,amssymb}
\usepackage{xcolor}
\usepackage{enumitem}
\usepackage[hidelinks]{hyperref}

% ---------- Eigene Makros ----------
\newcommand{\inc}{\texttt{inc}}
\newcommand{\bit}{\texttt{b1typ}}
\newcommand{\bzt}{\texttt{b2typ}}
\newcommand{\bic}{\texttt{b1cat}}
\newcommand{\bzc}{\texttt{b2cat}}
\newcommand{\cat}{\texttt{cat}}
\newcommand{\typ}{\texttt{typ}}
\newcommand{\ok}{\textcolor{green!55!black}{$\checkmark$}}
\newcommand{\nok}{\textcolor{red!75!black}{$\times$}}

% ================================================================
\begin{document}

\begin{titlepage}
\centering
\vspace*{2cm}
{\Large\bfseries Analyseprotokoll\par}
\vspace{0.5cm}
{\LARGE\bfseries Q-Learning Lauf 1\par}
\vspace{0.4cm}
{\large Robotergestützte Sortierung auf zwei Puffer\par}
\vspace{2.5cm}

\begin{tabularx}{0.9\textwidth}{@{}lX@{}}
\toprule
\textbf{Gegenstand}      & Auswertung der Q-Tabelle des ersten tabellarischen
                           Q-Learning-Trainings \\[2pt]
\textbf{Anwendungsfall}  & Sortierung eingehender Teile auf zwei Puffer,
                           Fehlläufer über Return-Loop \\[2pt]
\textbf{Datenbasis}      & Q-Tabelle als CSV, 432 Zustände $\times$ 3 Aktionen \\[2pt]
\textbf{Zustandsraum}    & $(\inc,\bit,\bzt,\bic,\bzc)$, erzeugt über
                           verschachtelte Schleifen ($\inc$ außen, $\bzc$ innen) \\[2pt]
\textbf{Aktionen}        & $A_0$ = Puffer 1, $A_1$ = Puffer 2, $A_2$ = Return-Loop \\[2pt]
\textbf{Methode}         & Dekodierung des Zeilenindex, Soll-Ist-Vergleich der
                           $\arg\max$-Policy gegen fachliche Regel,
                           Monotonieprüfung der Wertefunktion \\[2pt]
\textbf{Datum}           & 21.\,August 2026 \\
\bottomrule
\end{tabularx}

\vfill
\end{titlepage}

\tableofcontents
\clearpage

% ================================================================
\section{Ausgangslage und Zielsetzung}

Ziel des Trainings war es, ohne vorgegebene Heuristik eine Zuweisungsstrategie zu
erlernen, die eingehende Teile so auf zwei Puffer verteilt, dass möglichst viele
vollständige 10er-Auslagerungen entstehen und die Return-Loop möglichst selten
genutzt wird.

Geprüft wurden drei Fragen:
\begin{enumerate}[label=\arabic*.,leftmargin=*]
  \item Hat der Agent eine fachlich sinnvolle Strategie gelernt?
  \item Nutzt er die Diskretisierung der Füllstände ($\cat$) tatsächlich?
  \item Welche Zustände sind unzureichend gelernt und warum?
\end{enumerate}

% ================================================================
\section{Abdeckung des Zustandsraums}

\begin{table}[htbp]
\centering
\caption{Aufteilung der 432 Tabellenzeilen nach Erreichbarkeit und Besuchsstatus}
\begin{tabularx}{\textwidth}{@{}Xrrrl@{}}
\toprule
\textbf{Kategorie} & \textbf{Zust.} & \textbf{Anteil} &
\textbf{Besucht} & \textbf{Bewertung} \\
\midrule
Strukturell unmöglich ($\typ=0 \wedge \cat>1$) & 132 & 30{,}6\,\% & 0
  & korrekt, exakt $0{,}0$ \\
Betriebsrelevant (Puffer leer / ungl.\ Typen)   & 219 & 50{,}7\,\% & 219 (100\,\%)
  & vollständig gelernt \\
Beide Puffer gleicher Typ                       &  81 & 18{,}8\,\% & 49 (60\,\%)
  & $Q<66$, meist $<3$ \\
\quad davon nie besucht                         &  32 &  7{,}4\,\% & 0
  & --- \\
\bottomrule
\end{tabularx}
\end{table}

\paragraph{Feststellung.}
Der betriebsrelevante Teilraum ist zu 100\,\% erschlossen. Die Zustände
\enquote{beide Puffer enthalten denselben Typ} treten unter der gelernten Policy
praktisch nicht auf -- der Agent hat implizit gelernt, dass das Aufteilen eines
Typs auf zwei Puffer beide Auslagerungen verzögert. Die geringe Abdeckung dieser
81 Zustände ist damit ein \emph{Resultat der Policy, kein Trainingsmangel}.

% ================================================================
\section{Was der Agent gelernt hat}

\subsection{Prioritätskette der Zuweisung}

\begin{table}[htbp]
\centering
\caption{Soll-Ist-Vergleich der $\arg\max$-Policy über alle 216 trainierten Zustände}
\begin{tabularx}{\textwidth}{@{}XL{3.6cm}rr@{}}
\toprule
\textbf{Situation} & \textbf{Korrekte Aktion} &
\textbf{Treffer} & \textbf{Quote} \\
\midrule
$\inc = \bit$ (Match Puffer 1)      & Puffer 1            & 61 / 63 & \textbf{96{,}8\,\%} \\
$\inc = \bzt$ (Match Puffer 2)      & Puffer 2            & 60 / 63 & \textbf{95{,}2\,\%} \\
ein Puffer leer, anderer Fremdtyp   & leeren Puffer belegen & 33 / 36 & \textbf{91{,}7\,\%} \\
beide belegt, kein Match            & Return              & 37 / 54 & 68{,}5\,\% \\
beide leer                          & symmetrisch beliebig &  3 / 3  & --- \\
\midrule
\textbf{Gesamt}                     &                     & \textbf{194 / 216} & \textbf{88{,}4\,\%} \\
\bottomrule
\end{tabularx}
\end{table}

Der Agent hat die Regel \textbf{Match $\rightarrow$ freier Puffer $\rightarrow$
Return} selbstständig und quasi-deterministisch gefunden.

\subsection{Nichttriviales Zusatzverhalten: gelerntes Opfer-Kriterium}

Die 17 Abweichungen in der Kategorie \enquote{beide belegt, kein Match} sind
überwiegend keine Fehler. In 9 von 11 Fällen mit ungleichem Füllstand
überschreibt der Agent gezielt den Puffer mit dem \emph{niedrigeren}
$\cat$-Wert, statt die Return-Loop zu nutzen.

\begin{table}[htbp]
\centering
\caption{Zustände mit gelerntem Opfer-Kriterium (fett: niedrigerer Füllstand)}
\begin{tabular}{@{}rlll@{}}
\toprule
\textbf{Zeile} & \textbf{Zustand $(\inc,\bit,\bzt,\bic,\bzc)$} &
\textbf{Aktion} & \textbf{Interpretation} \\
\midrule
101 & 1,\,2,\,3,\,\textbf{1},\,3 & Puffer 1 & opfert den fast leeren Puffer \\
106 & 1,\,2,\,3,\,3,\,\textbf{1} & Puffer 2 & opfert den fast leeren Puffer \\
128 & 1,\,3,\,2,\,\textbf{1},\,3 & Puffer 1 & dito \\
129 & 1,\,3,\,2,\,2,\,\textbf{1} & Puffer 2 & dito \\
212 & 2,\,\dots,\,\textbf{2},\,3 & Puffer 1 & dito \\
347 & 3,\,\dots,\,\textbf{2},\,3 & Puffer 1 & dito \\
348 & 3,\,1,\,2,\,3,\,\textbf{1} & Puffer 2 & dito \\
375 & 3,\,2,\,1,\,3,\,\textbf{1} & Puffer 2 & dito \\
\bottomrule
\end{tabular}
\end{table}

Der Agent wägt also \textbf{Rückführkosten gegen Sortierfortschritt} ab -- ein
Verhalten, das eine handgeschriebene \enquote{Match-first / sonst Return}-Heuristik
nicht enthält. Werden diese 9 Fälle als sinnvoll gewertet, steigt die
Policy-Qualität auf \textbf{92{,}6\,\%}.

\subsection{Wertefunktion nutzt die Füllstandsdiskretisierung}

\begin{table}[htbp]
\centering
\caption{$Q$-Wert der gewählten Aktion für $\inc=1$, $\bit=2$, $\bzt=1$
         (Match auf Puffer 2, Zeilen 81--89)}
\begin{tabular}{@{}cccc@{}}
\toprule
\multirow{2}{*}{$\bic$} & \multicolumn{3}{c}{$\bzc$} \\
\cmidrule(l){2-4}
 & \textbf{1} & \textbf{2} & \textbf{3} \\
\midrule
\textbf{1} & 605{,}3 & 625{,}4 & 659{,}6 \\
\textbf{2} & 625{,}8 & 645{,}0 & 669{,}2 \\
\textbf{3} & 641{,}7 & 669{,}2 & \textbf{690{,}1} \\
\bottomrule
\end{tabular}
\end{table}

Die Werte steigen monoton in \emph{beiden} Füllständen. Daraus folgt:
\begin{itemize}[leftmargin=*]
  \item Der Füllstand des Zielpuffers wird korrekt als Nähe zur
        10er-Auslagerung bewertet.
  \item Auch der Füllstand des \emph{anderen} Puffers erhöht den Zustandswert
        $\Rightarrow$ der Agent besitzt einen Fortschrittsbegriff für den
        Gesamtzustand, nicht nur eine lokale Regel.
  \item \textbf{Nachweis, dass die $\cat$-Features informativ genutzt werden}
        und kein toter Ballast im Zustandsvektor sind.
\end{itemize}

Wertespanne der trainierten Zustände: $509{,}5$ (Zeile 180) bis $700{,}2$
(Zeile 125). Return-$Q$-Werte liegen mit $511$--$687$ im gleichen Band; die
Schleife wird als moderat teuer, nicht als katastrophal bewertet.

% ================================================================
\section{Was der Agent nicht gelernt hat}

\subsection{Fehlentscheidungen -- alle mit identischer Ursache}
\label{sec:fehler}

\begin{table}[htbp]
\centering
\caption{Alle 8 Fehlentscheidungen der gelernten Policy}
\begin{tabular}{@{}rlcrr@{}}
\toprule
\textbf{Zeile} & \textbf{Zustand $(\inc,\bit,\bzt,\bic,\bzc)$} &
\textbf{Gewählt} & \textbf{$Q$ gewählt} & \textbf{$Q$ korrekt} \\
\midrule
 29 & 1,\,0,\,3,\,1,\,\textbf{3} & Return   & 606{,}3 & P1 = \textbf{0{,}00} \\
173 & 2,\,0,\,3,\,1,\,\textbf{3} & Return   & 612{,}1 & P1 = 0{,}26 \\
299 & 3,\,0,\,1,\,1,\,\textbf{3} & Return   & 627{,}3 & P1 = 0{,}05 \\
245 & 2,\,2,\,3,\,1,\,\textbf{3} & Return   & 624{,}7 & P1 = 0{,}73 \\
 61 & 1,\,1,\,2,\,\textbf{3},\,2 & Puffer 2 & 544{,}0 & P1 = 2{,}14 \\
353 & 3,\,1,\,3,\,1,\,\textbf{3} & Puffer 1 & 611{,}0 & P2 = 9{,}59 \\
356 & 3,\,1,\,3,\,2,\,\textbf{3} & Puffer 1 & 607{,}7 & P2 = 2{,}05 \\
359 & 3,\,1,\,3,\,3,\,\textbf{3} & Puffer 1 & 598{,}1 & P2 = 13{,}50 \\
\bottomrule
\end{tabular}
\end{table}

\end{document}
